\ssszh 处理完必要的修正，现在我们终于兑现承诺，来实际看看图 \ref{f:table} 中各对象之间缠绕算子构成的范畴。这样做，我们会遇到许多漂亮的定理与猜想，希望这能让读者对本学科有更多亲自动手的感觉。

两侧出现的范畴都是自对偶的。对 D-模与拟凝聚层的范畴，除开一些可驾驭的无穷维问题，这由例 \ref{ex:dmodselfdual} 得出。Kac--Moody 表示的情形我们会在下文第 \ref{ss:affhc} 节更详细地讨论。承认这种自对偶，经由典范等价
\[
\Hom_{\Dmod(LG)\mod}(\sC, \sD) \simeq \Hom_{\Dmod(LG)\mod}(\sD^\vee, \sC^\vee),
\]
以及对 $\LS$ 上的范畴类似的关系，我们只需写下「对角线上方」的缠绕算子范畴。

这些缠绕算子记录在图 \ref{f:tableaut} 与图 \ref{f:tablespec} 中。自守侧与谱侧这些范畴的计算，基本上遵循第 \ref{s:catrepgrp} 节、尤其是例 \ref{e:fibprodvar} 与第 \ref{ss:invariantscosets} 节讨论的模式。下面我们来比较两张表中的条目，并对出现的等价（无论已知还是未知）加以评注。

\begin{table}[H]
\centering
\scriptsize
\setlength{\tabcolsep}{1.6pt}
\begin{tabular}{|c|c|c|c|c|}
\hline
 & $\Dmod(\Gr_G)$ & $\Dmod(\Fl_G)$ & $\Dmod(LG/LN, \psi)$ & $\fgc\mod$ \\
\hline
$\Dmod(\Gr_G)$ & $\Dmod(L^+G \backslash LG / L^+G)$ & $\Dmod(L^+G
\backslash LG / I)$ & $\Dmod(L^+G \backslash LG / LN, \psi)$ &
$\fgc\mod^{L^+G}$ \\
\hline
$\Dmod(\Fl_G)$ & & $\Dmod(I \backslash LG / I)$ &
$\Dmod(I \backslash LG / LN, \psi)$ & $\fgc\mod^I$ \\
\hline
$\Dmod(LG/LN, \psi)$ & & & $\Dmod(LN, \psi \backslash
LG / LN, \psi)$ & $\fgc\mod^{LN, \psi}$ \\
\hline
$\fgc\mod$ & & & & $\fgc\oplus
\fgc\mod^{LG}$ \\
\hline
\end{tabular}
\caption{局部几何朗兰兹自守侧的一些缠绕算子。}
\label{f:tableaut}
\end{table}

\begin{table}[H]
\centering
\scriptsize
\setlength{\tabcolsep}{1.6pt}
\begin{tabular}{|c|c|c|c|c|}
\hline
 & $\QCoh(\on{pt}/\cG)$ & $\QCoh(\widetilde{\sN} / \cG)$
& $\QCoh(\LS)$ & $\QCoh(\Op)$ \\
\hline
$\QCoh(\on{pt} / \cG)$ & $\ICn( \on{pt} / \cG \underset {\cfg/\cG} \times
\on{pt} / \cG)$ & $\ICn(\on{pt} / \cG \underset { \cfg / \cG} \times
\widetilde{\sN} / \cG)$ & $\QCoh( \on{pt}/ \cG)$ &
$\QCoh( \Op^{\on{unramified}})$ \\
\hline
$\QCoh(\widetilde{\sN} / \cG)$ & & $\ICn( \widetilde{\sN}/ \cG \underset{
\cfg / \cG } \times \widetilde{\sN} / \cG)$ & $\QCoh( \widetilde{\sN}/
\cG)$ & $\QCoh(\widetilde{\sN}/\cG \underset \LS \times \Op)$ \\
\hline
$\QCoh(\LS)$ & & & $\QCoh(\LS)$ & $\QCoh(\Op)$ \\
\hline
$\QCoh(\Op)$ & & & & $\QCoh( \Op \underset \LS \times \Op)$ \\
\hline
\end{tabular}
\caption{局部几何朗兰兹谱侧的一些缠绕算子。}
\label{f:tablespec}
\end{table}

\ssszh 从非分歧情形开始，即表中的第一行。第一个等价
\begin{equation}
\Dmod(L^+G \backslash LG / L^+G) \simeq \ICn( \on{pt} / \cG \underset {\cfg/\cG}
\times \on{pt} / \cG)
\end{equation}
我们已经遇到过，即 Bezrukavnikov--Finkelberg \cite{bezfinksat} 的导出 Satake 等价。

导出 Satake 等价是 $t$-正合的，在阿贝尔范畴上恢复几何 Satake 等价
\[
 \Dmod(L^+G \backslash LG / L^+G)^\heartsuit \simeq \QCoh(\on{pt} / \cG) =
 \on{Rep}(\cG),
\]
后者是更早的一个基本定理，属于 Lusztig \cite{lusztigqchar}、Drinfeld、Ginzburg \cite{ginzburgsatake}、Mirković--Vilonen \cite{mirkovic-vilonen}。这两个等价（其中阿贝尔的那个稍更透明）都是 $p$-adic 情形中类似断言的范畴化：
\[
   \on{Dist}( G(\Fq[[t]]) \backslash G(\Fq(\!(t)\!)) / G(\Fq[[t]])) \simeq K_0(
   \on{Rep}(\cG)) \underset{\mathbb{Z}}{\otimes} k,
\]
该同构连同其混合特征的情形，属于 Satake \cite{satake63}。

对此作两点评注。

\begin{re} \label{r:paramcount} 第一，对初次接触的读者，为了对这个等价有所把握，配比参数（即（几乎）紧对象的 Grothendieck 群的基）是很有启发的。

基本要点是：$LG$ 上只有可数多个 $L^+G \times L^+G$-轨道，它们自然地以 $\cG$ 的不可约表示为指标。更仔细地说，任何余特征标
\[
   \check{\mu}: \mathbb{G}_m \rightarrow T
\]
限制到原点处的形式圆盘，就给出 $LT$ 的一点，记作 $t^{\check{\mu}}$。若以 $\check{\Lambda}^+$ 记支配余特征标，则 $LG$ 被轨道
\[
   L^+G \cdot t^{\check{\mu}} \cdot L^+G, \quad \text{对 } \check{\mu} \in
   \check{\Lambda}^+,
\]
分层。而这些正是 $\cG$-模的最高权；该等价把 $\Gr_G$ 上相应的相交上同调复形与单 $\cG$-模互换。\end{re}

\begin{re} 第二，这是一条帮助定位的注。导出 Satake 自守侧底层的空间，参数化三元组 $(\sP_1, \sP_2, \tau)$，其中 $\sP_1, \sP_2$ 是圆盘上的 $G$-丛，$\tau$ 是它们在穿孔圆盘上的等同。

类似地，谱侧底层的空间参数化三元组 $(\sE_1, \sE_2, \tau)$，其中 $\sE_1, \sE_2$ 是形式圆盘上的 $\cG$-联络，$\tau$ 是它们限制到穿孔圆盘之后的同构。注意，由于形式圆盘上每个 $\cG$-联络都可平凡化，底层的经典叠是 $\on{pt} / \cG$。然而这个模问题要求我们把 $\on{pt}/\cG$ 与它在 $\LS$ 中的自身相交，正是这提供了使该等价成立的导出结构。

因此，不正式地说，这个等价把 $G$ 的一个凝聚模空间上的构造性对象，与 $\cG$ 的类似构造性模空间上的凝聚对象互换。
\end{re}

\ssszh \label{s:abg}表中的下一个等价，即
\[
   \Dmod(L^+G \backslash LG / I) \simeq \ICn(\on{pt}/\cG \underset {\cfg / \cG}
   \times \widetilde{\sN} / \cG),
\]
也是 Arkhipov--Bezrukavnikov--Ginzburg \cite{abg} 一个非常好的定理。如注 \ref{r:paramcount} 那样，可以这样配比参数。在自守侧，$LG$ 的相关分层现在由层
\[
     L^+ G \cdot t^{\check{\mu}} \cdot I, \quad \text{对 } \check{\mu} \in
     \check{\Lambda},
\]
给出。也就是说，要把右侧的 $L^+G$ 缩小到 $I$，就必须乘以有限 Weyl 群这个因子，即从支配余特征标扩大到全部余特征标。

在谱侧，回忆 $\widetilde{\sN}$ 是（有限）旗簇的余切丛，我们考虑的是零截面（即旗簇本身）的导出加厚上的 $\cG$-等变凝聚层。这里，若忽略加厚，同一个格 $\check{\Lambda}$ 出现在
\[
   \QCoh( \Fl_{\cG}  / \cG  ) \simeq \QCoh(  \cB \backslash \cG / \cG) \simeq
   \QCoh( \cB \backslash \on{pt}) \simeq \on{Rep}(\cB),
\]
中，即作为 $\Fl_{\cG}$ 上的等变线丛，等价地作为 $\on{Rep}(\cB)$ 的特征标。
